\chapter*{前言}
\chapterauthor{國立臺灣師範大學 資訊工程學系 紀博文}
\label{chap:forward}
\addcontentsline{toc}{chapter}{前言}

第一次聽到這個課名，是在執行教育部資訊科技第二專長學分班的計畫時聽到的。在台師大的「國立臺灣師範大學中等學校師資職前教育專門課程『科技領域資訊科技專長』」裡面規定有一門領域核心必修課程：「科技系統與社會發展」，但我對於這門課的內容卻一無所知，自然也不知道找哪個老師來協這開設這門課。但計畫還是要執行的，所以我只好硬著頭皮去看科技領域的「十二年國民基本教育課程綱要」，然後我看到了下面這段文字：

\begin{quote}
資訊科技課程亦協助學生建立資訊社會中應有的態度，透過對資訊科技與人類社會相關議題之了解，養成正確的資訊科技使用習慣，遵守相關之倫理、道德及法律，並關懷資訊社會的各項議題。
\end{quote}

看起來，這門課應該是屬於社會學領域的「科技與社會（英語：Science, technology and society, STS）」議題。所以我就很高興的去看看其他學校的相關課程怎麼開，想說依樣畫葫蘆即可。結果發現，我大概是沒有慧根，完全看不懂課綱在寫什麼，一堆社會學的專用術語，像是「問題意識」、「社會層次」、「性別政治」等，對於工程頭腦的我完全無法理解。最終我決定把收集到的資料丟到垃圾桶，重新去反思對於理科、工程腦袋的學生來說，這門課應有的長相。

\section{科技真的那麼美好？}

身為學術界的一員，發表論文是不得不做的事情。而在發表論文的時候，通常都會寫一段「廣告台詞」，試圖說服讀者說這個技術是非常重要的，有了這個技術，就可以解決人類社會的關鍵難題，帶來無比美好的未來。但這是真的嗎？？

我永遠記得在唸書時曾經修過一門「物聯網概論」，老師在第一堂課就說：「物聯網技術的誕生可以大幅提昇人類生活的便利性，帶給人更輕鬆的生活。」聽起來很有道理，但物聯網逐漸成熟的現在，美好的未來真的發生了嗎？看起來似乎沒有，自動化取代了許多簡單的工作，迫使人們承擔更多複雜的任務，並且步調也越來越快，壓力越來越大。如果從這個角度出發，物聯網技術真的值得研究嗎？

所以在這門課，我希望從各種面向去討論科技對於人類社會的各種影響，帶這些理科、工程腦袋的學生去思考技術對於真實世界的衝擊。我盡可能站在學生的對立面，去挑戰學生們的立場跟想法，強迫讓他們去反思技術對於人類社會真的是好事嗎？

\section{一隻手的經濟學家}

美國總統 Harry Truman 曾說過：

\begin{quote}
Give me a one-handed Economist. All my economists say 'on ONE hand...', then 'but on the other...
\end{quote}

就我的理解，這句話的意思是經濟學家往往喜歡把正反雙方的意見並陳，然後打死不做最後的決定，同樣的情況也出現在哲學、神學書籍，這些作者往往介紹了各種理論的內容，但對於衝突卻往往用「張力」兩個字帶過。這種四平八穩的回覆非常官方，但卻給人一種不負責任的感覺。

因此，在這門課裡面，我要求學生們要針對特定科技對社會的衝擊進行多方思考，但最終要做出自己的選擇。我不要求他們要說服我，畢竟一百個人可能有一百零一種看法，學生的看法跟不我同是很自然的事情。我要的是他們知不知道自己選擇所帶來的優點以及附上的代價，讓他們自己察覺，有得也一定會有失。這種 trade-off 的概念也是一個工程師該有的素養。

\section{文章架構}

文章的基本結構如下。首先，會請學生介紹技術對社會帶來的衝擊，這邊會要求他們對正反雙方（或是更多面向）進行簡單的說明。然後會要求他們針對技術做科普等級的介紹，畢竟任何的討論都必須建立在對技術正確的認識上，對技術的錯誤認識可能會導出莫名其妙的結論。接下來會要求他們深入檢視各方的想法，最終提出自己的結論。

為了讓同學不會活在自己的世界當中，我希望他們能夠進可能的跟這個世界對話。所以每篇文章我會幫同學找各領域（不見得是電機資訊）的專家，讓他們對文章來給出意見。這裡不是做學生的報告進行評分，而是對話、討論，針對學生的意見找出邏輯上的漏洞或是提出不一樣論點，最後會在要求學生進行回應，希望這樣能夠刺激學生去思考自己的選擇有沒有調整的可能。

本書就是在這個概念下誕生的成果，但願我有足夠的恆心毅力把這個課程持續進行下去。

